% Additive extension: action, area, sectorial memory, and thermometry.
% Owners: 70_horizontes_e_informacion.tex; preceding thermal section;
% 71_energia_libre_holografica_cerrada.tex in the integral corpus.
% Conditions of the declared cosmological realization are retained.

\subsection{Thermal conversion between the calendar and the horizon}
\label{viii:subsec:conversion-termica-calendario-horizonte}

Action, the area increment, and angular memory admit an explicit
energetic composition. We retain one action orientation, the positive
Barbero--Immirzi branch, and the realization of
\eqref{vii:eq:horizonte-cosmologico}. Write
\[
 L=\ell_P=L_\alpha,\qquad
 E_P=\frac{h}{108t_0}=\frac{\hbar c}{L},\qquad
 q_{\mathcal A}=\delta\mathcal A=\gamma_{\rm BI}a_0L^2.
\]
The symbol \(q_{\mathcal A}\) denotes an area increment here.
The angular factor \(a_0\), the functional \(\gamma_{\rm BI}\), and
action are the quantities obtained in the preceding sections.
Their domain preserves the hexadic--octadic incidence, orientation,
and calendar of the same enriched APP--TRIT--TPK state.

\begin{proposition}[Composition of the areal and thermal scales]
\label{viii:prop:composicion-area-termica}
For the tension in~\eqref{vii:eq:tension-informacion-area},
\begin{equation}
 \frac{\mathcal T_{I/A}q_{\mathcal A}}{\log3}
 =\frac{E_P}{\log3}
 =\frac{h}{108t_0\log3}
 =k_B\Theta_{\rm clk}.
 \label{viii:eq:composicion-area-termica}
\end{equation}
For \(R=\zeta L\), with \(\zeta>0\), let
\(\tau_R=k_BT_{\rm cos}(R)\). Then
\begin{equation}
 \tau_R=\frac{E_P}{2\pi\zeta},\qquad
 \frac{T_{\rm cos}(R)}{\Theta_{\rm clk}}
 =\frac{\log3}{2\pi\zeta}.
 \label{viii:eq:conversion-dos-temperaturas}
\end{equation}
In particular, the ratio at \(R=L\) is \(\log3/(2\pi)\).
\end{proposition}
\begin{proof}
The definition \(\mathcal T_{I/A}=E_P/q_{\mathcal A}\) and the
tritic normalization in~\eqref{vii:eq:unidad-accion-informacion}
give~\eqref{viii:eq:composicion-area-termica}. The gravitational
expression of the same tension gives the same result:
\[
 \frac{c^4}{\gamma_{\rm BI}a_0GL}
        \gamma_{\rm BI}a_0L^2
 =\frac{c^4L}{G}=\frac{\hbar c}{L}=E_P,
\]
using \(G\hbar=c^3L^2\).
By~\eqref{vii:eq:horizonte-cosmologico},
\(\tau_R=\hbar c/(2\pi R)=E_P/(2\pi\zeta)\).
Dividing this identity by \(k_B\Theta_{\rm clk}=E_P/\log3\)
proves the temperature ratio.
\end{proof}

In the same action section, the gravitational composition reads
\begin{equation}
 Gk_B\Theta_{\rm clk}\log3=GE_P=c^4L,
 \qquad G\tau_R=\frac{c^4L}{2\pi\zeta}.
 \label{viii:eq:composicion-termica-gravitatoria}
\end{equation}
Both identities follow from \(G\hbar=c^3L^2\) and the preceding
scales.

The tension converts the area increment into decision energy.
Tritic normalization and cosmological periodicity determine two thermal
sections of that same energy, related by
\eqref{viii:eq:conversion-dos-temperaturas}. The conversion preserves
both normalizations and the role of Barbero--Immirzi in geometric
resolution.

\subsection{Energetic response of sectorial memory}
\label{viii:subsec:respuesta-energetica-sectorial}

Fix \(R>0\) and the faithful state
\(\rho_0=\rho_\partial(\Delta_0)\), with \(\Delta_0>0\),
in~\eqref{vii:eq:estado-sectorial}. Its modular operator is
\(\mathsf K_0=(\log Z_0)I+\Delta_0\mathsf D_\partial\).
The horizon realization assigns to each density state
\(\sigma\) the response quantities
\begin{equation}
\begin{aligned}
 \mathcal E_R(\sigma)
 &=E_\Lambda(R)+\tau_R
   \operatorname{Tr}[(\sigma-\rho_0)\mathsf K_0],\\
 \mathcal S_R(\sigma)
 &=S_{\rm H}(R)+k_B[s(\sigma)-s(\rho_0)].
\end{aligned}
 \label{viii:eq:respuesta-horizonte-definida}
\end{equation}
These describe the response of the modular state in the geometric
realization already specified. Radius and action fix its energetic
scale \(\tau_R\); occupations preserve the provenance of each channel.

\Needspace{28\baselineskip}
\begin{theorem}[Response Hamiltonian and informational balance]
\label{viii:thm:hamiltoniano-respuesta-termica}
The operator of energy differences corresponding to
\eqref{viii:eq:respuesta-horizonte-definida} may be chosen as
\begin{equation}
 \mathsf H_R=\tau_R\Delta_0\mathsf D_\partial.
 \label{viii:eq:hamiltoniano-respuesta}
\end{equation}
Its Gibbs state at energy scale \(\tau_R\) is \(\rho_0\).
For every density state \(\sigma\) on the sectorial space,
\begin{equation}
 \mathcal E_R(\sigma)-T_{\rm cos}(R)\mathcal S_R(\sigma)
 =\tau_R\mathcal D(\sigma\Vert\rho_0)\geq0.
 \label{viii:eq:balance-libre-termico}
\end{equation}
For a differentiable state variation about \(\rho_0\), with
\(R\) and \(\mathsf H_R\) fixed, one obtains
\begin{equation}
 \delta\mathcal E_R
 =\tau_R\,\delta s
 =T_{\rm cos}(R)\,\delta\mathcal S_R.
 \label{viii:eq:primera-ley-energetica}
\end{equation}
\end{theorem}
\begin{proof}
The traces of \(\sigma\) and \(\rho_0\) are one, so their
difference eliminates the scalar term \(\log Z_0\). Hence
\(\mathcal E_R(\sigma)-E_\Lambda(R)
=\operatorname{Tr}[(\sigma-\rho_0)\mathsf H_R]\).
Moreover,
\[
 \frac{e^{-\mathsf H_R/\tau_R}}
 {\operatorname{Tr}e^{-\mathsf H_R/\tau_R}}
 =\frac{e^{-\Delta_0\mathsf D_\partial}}{Z_0}=\rho_0.
\]
The identity \(E_\Lambda=T_{\rm cos}S_{\rm H}\) cancels the
geometric reference in the balance. The finite identity
\eqref{vii:eq:ley-modular-finita}, written as
\(s(\sigma)-s(\rho_0)
=\operatorname{Tr}[(\sigma-\rho_0)\mathsf K_0]
-\mathcal D(\sigma\Vert\rho_0)\), proves
\eqref{viii:eq:balance-libre-termico}. Its tangent term is the first
modular law already established and gives
\eqref{viii:eq:primera-ley-energetica}.
\end{proof}

For a unit occupation increment, the energy gaps are
\(\epsilon_{90}=90\tau_R\Delta_0\) and
\(\epsilon_{120}=120\tau_R\Delta_0\). In both sectors,
\begin{equation}
 \frac{\epsilon_m}{m\Delta_0}=\tau_R,
 \qquad m\in\{90,120\}.
 \label{viii:eq:pendiente-termica-comun}
\end{equation}
Thus the ratio \(4/3\) between areal weights records different gaps
of the same Hamiltonian. The ratio of an energy gap to the corresponding
variation of \(-\log p\) preserves a common scale. State normalization
eliminates only a scalar constant.

\begin{corollary}[Energy susceptibility at fixed Hamiltonian]
\label{viii:cor:respuesta-escala-termica}
For \(\rho_\tau=e^{-\mathsf H_R/\tau}/
\operatorname{Tr}e^{-\mathsf H_R/\tau}\), with \(\tau>0\),
\begin{equation}
 \frac{\mathrm d\langle\mathsf H_R\rangle_\tau}{\mathrm d\tau}
 =\frac{\operatorname{Var}_{\rho_\tau}(\mathsf H_R)}{\tau^2}>0.
 \label{viii:eq:susceptibilidad-termica}
\end{equation}
At \(\tau=\tau_R\), the result is
\(\Delta_0^2\operatorname{Var}_{\rho_0}(\mathsf D_\partial)\).
\end{corollary}
\begin{proof}
Differentiating the finite sum of exponentials gives the variance
divided by \(\tau^2\). The state is faithful and \(\mathsf H_R\)
has more than one eigenvalue, so the variance is positive. Substitution
of~\eqref{viii:eq:hamiltoniano-respuesta} gives the stated evaluation.
\end{proof}

This derivative varies the Gibbs scale while holding the Hamiltonian
fixed. The radial variation in~\eqref{vii:eq:diferencial-horizonte}
also changes \(\mathsf H_R\) and retains its own geometric terms.

\subsection{Thermometric section and preservation of the record}
\label{viii:subsec:seccion-termometrica-registro}

The state \(\rho_0\) and response Hamiltonian determine a thermometric
coordinate before any candidate numerical reader is considered.
For the fixed radius, let \(\mathcal L_\Theta=\mathbb R u_R\) be
an oriented line whose positive section \(u_R\) corresponds to the
reference state. On faithful Gibbs states of the same Hamiltonian,
\begin{equation}
 \rho_\beta=\frac{e^{-\beta\mathsf H_R}}
                  {\operatorname{Tr}e^{-\beta\mathsf H_R}},\qquad
 \Theta_R(\rho_\beta)=\frac{1}{\beta\tau_R}u_R,
 \quad \beta>0.
 \label{viii:eq:lector-termometrico}
\end{equation}
The operator identity
\(-\log\rho_\beta=\beta\mathsf H_R+(\log Z_\beta)I\)
determines \(\beta\): subtracting two eigenvalues with distinct
energies gives its unique ratio. In particular,
\(\Theta_R(\rho_0)=u_R\).
The line \(\mathcal L_E\) is the one-dimensional space of energetic
quantities.

\Needspace{23\baselineskip}
\begin{theorem}[Positive transducer and identification criterion]
\label{viii:thm:identificacion-transductor}
There is a unique positive linear map
\(B_R:\mathcal L_\Theta\to\mathcal L_E\) satisfying
\(B_R(u_R)=\tau_R\). It obeys
\begin{equation}
 B_R(\Theta_R(\rho_\beta))=\beta^{-1}.
 \label{viii:eq:transductor-gibbs}
\end{equation}
The compatible chronological section is
\begin{equation}
 \Theta_{{\rm clk},R}
 =\frac{E_P}{\tau_R\log3}u_R
 =\frac{2\pi R}{L\log3}u_R,
 \qquad
 B_R(\Theta_{{\rm clk},R})=
 \frac{\mathcal T_{I/A}q_{\mathcal A}}{\log3}.
 \label{viii:eq:seccion-cronologica-compatible}
\end{equation}
For any other linear map \(B_*\) between the same lines,
\begin{equation}
 B_*=B_R
 \quad\Longleftrightarrow\quad
 B_*(\Theta_{{\rm clk},R})=\frac{h}{108t_0\log3}.
 \label{viii:eq:criterio-una-seccion}
\end{equation}
\end{theorem}
\begin{proof}
The section \(u_R\) is a basis of a line, so prescribing its
image defines a unique linear map; positivity of \(\tau_R\)
ensures positivity of the map. Substituting
\eqref{viii:eq:lector-termometrico} proves
\eqref{viii:eq:transductor-gibbs}. The expression for \(\tau_R\)
and~\eqref{viii:eq:composicion-area-termica} give
\eqref{viii:eq:seccion-cronologica-compatible}. Finally,
\(\Theta_{{\rm clk},R}\) is nonzero. Two linear maps from a line
coincide if and only if they coincide on that section.
\end{proof}

The criterion compares readers on a section previously obtained from
the state, energetic response, and tritic normalization. Application to
a concrete numerical coordinate requires evaluating that reader on
the specified section. Dimensional bases are transported jointly with
the maps and the thermal coordinates they express.

\begin{proposition}[Transport of the response with memory]
\label{viii:prop:transporte-termico-memoria}
Let \(J\) be an isometry preserving the complete record labels,
and take its image as the receiving space. For
\(\mathsf H'_R=J\mathsf H_RJ^*\), \(\rho'_0=J\rho_0J^*\),
and \(\sigma'=J\sigma J^*\), energy expectations, entropies, and
balance~\eqref{viii:eq:balance-libre-termico} are preserved.
Moreover, \(J\mathsf H_R=\mathsf H'_RJ\).
\end{proposition}
\begin{proof}
On its image, \(J\) is unitary and transports functional calculus.
Cyclicity of the trace preserves expectations; the same nonzero
eigenvalues preserve entropy and relative entropy. The intertwining
identity uses \(J^*J=I\). Substitution into the balance proves its
preservation.
\end{proof}

The same equivalence admits another scale \(\tau'=a\tau_R\),
with \(a>0\), and another energy origin \(e'\). If
\(\mathsf H'=e'I+aJ\mathsf H_RJ^*\), then
\[
 \frac{\mathsf H'-e'I}{\tau'}
 =J\frac{\mathsf H_R}{\tau_R}J^*,\qquad
 \frac{e^{-\mathsf H'/\tau'}}
      {\operatorname{Tr}_{\operatorname{im}J}e^{-\mathsf H'/\tau'}}
 =J\rho_0J^*.
\]
The scalar factor \(e^{-e'/\tau'}\) cancels in the state;
energy differences are multiplied by \(a\), while relative entropy
remains invariant. The corresponding balance therefore has coefficient
\(\tau'\). The full energy expectation preserves the origin:
\(\langle\mathsf H'\rangle=e'+a\langle\mathsf H_R\rangle\).
Realizations are thus compared through their centered exponents and
their own scales.
If the expectation supplies the readings \(m=E/c^2\) and
\(r_g=GE/c^4\), keeping \(c,G\) fixed gives
\[
 m'=\frac{e'}{c^2}+am,\qquad
 r'_g=\frac{Ge'}{c^4}+ar_g.
\]
Statistical normalization and the transport of quantities therefore
retain distinct roles for the same recorded origin.

The composition connects the Barbero--Immirzi area increment, calendar
energy, and angular boundary distribution. Its transport retains
sectorial provenance and the record: phase return remains accompanied
by the memory advance of the enriched state.
